\section{A singular power-law field}
\label{c:sing-fields}

The rational field smooths the common tip before tilting the lines.  We now
keep a power-law singularity in the derivative and remove a small tip
region instead.  Two triangular-base cylinders cover the removed regions;
their cost is smaller than the quadratic saving from the tilted families.
The exponent of the power law makes both requirements compatible.
We give the construction at edge length two, where the cone walls have
slopes $\pm1$, and then state its exact edge-one form.

We use the regular tetrahedron of edge length $2$,
\[
 K_2=\{(x,y,\sqrt{2}\,t):0\le t\le1,
                 \ |x|\le1-t,\ |y|\le t\}.
\]
The physical metric in these coordinates is $dx^2+dy^2+2\,dt^2$.
The minimum projection area is
$A_{\min}(K_2)=4A_{\min}(K_1)=\sqrt2$, since these regular tetrahedra
are related by a similarity of ratio two.  Appendix~\ref{c:projection-appendix}
gives two independent proofs of this normalization.

\begin{theorem}\label{c:sing-theorem}
For every
\[
 1<p<(\sqrt{51}-5)/2,\qquad 1<\gamma<1/(2-p),
\]
all sufficiently small positive $c$ determine a finite cylinder cover
of the regular edge-$2$ tetrahedron $K_2$ above with total perpendicular
base area less than $\sqrt2/2$.  The cover consists of two finite
families with compact parallelogram bases and two cylinders with
compact triangular bases.  Before the arbitrarily small errors needed
for the two finite approximations, the area bound is
\[
 \frac{\sqrt2}{2}+2\sqrt2\,\kappa_p c^2+o(c^2),
 \qquad
 \kappa_p=\frac{(1/2)^4(2p^2+10p-13)}{30p}<0.
\]
The asymptotic statement fixes $p$ and $\gamma$ first.  Each finite
approximation is chosen after the positive value of $c$ is fixed.
\end{theorem}

Put $a=1/2$.  The two cylinders obtained by splitting at $t=a$,
with axes parallel to the $x$- and $y$-axes respectively, have total
base area $2\sqrt{2}\,a^2=\sqrt{2}/2$.  We perturb those directions.

\subsection{The field and the cutoff}

Fix parameters in the order
\begin{equation}
 1<p<\frac{\sqrt{51}-5}{2},\qquad
 1<\gamma<\frac1{2-p},\qquad
 d=c^\gamma\quad(c>0).
 \label{c:sing-parameters}
\end{equation}
The upper bound on $p$ makes the quadratic area coefficient negative.
The bounds on $\gamma$ make the derivative of the tilted field uniformly
small while keeping the tip area $o(c^2)$.
Only $c$ will tend to zero.  To obtain a field tangent to the cone walls
$b=\pm r$, take a stream function that vanishes on both walls.
Multiplying $b^2-r^2$ by $r^{p-2}$ makes the stream function homogeneous
of degree $p$, so the resulting field has degree $p-1$.  The normalization
below also gives $Q(b,a)=b$, which will cancel the first-order overlap terms.
For $r>0$, define
\begin{align}
 \psi(b,r)&=\frac12(r/a)^{p-2}(b^2-r^2),
 &V=(P,Q)&=(-\partial_r\psi,\partial_b\psi),\notag\\
 P(b,r)&=\frac a2(r/a)^{p-1}
       \left[p+(2-p)(b/r)^2\right],
 &Q(b,r)&=(r/a)^{p-2}b.
 \label{c:sing-field}
\end{align}
At the dividing level let
\begin{equation}
 P_a(b)=P(b,a)=\frac{ap}{2}+\frac{2-p}{2a}b^2,
 \qquad \delta(b)=b^2P_a(b).
 \label{c:sing-ceiling-polynomial}
\end{equation}
Thus $\delta\ge0$, $Q(b,a)=b$, and $Q=\pm P$ on $b=\pm r$.
Our compact label domain is
\begin{equation}
 D_c=\{(b,r):d\le r\le2,\ |b|\le r,
                  \ r\le a+c^2\delta(b)+c^{5/2}\}.
 \label{c:sing-domain}
\end{equation}

The field is smooth on $r>0$, and direct differentiation gives
\begin{equation}
 \operatorname{tr}DV=0,\qquad
 \det DV=(r/a)^{2p-4}\frac{p-1}{2}
               \left[(2-p)(b/r)^2-p\right]<0
 \quad (|b|\le r).
 \label{c:sing-hyperbolicity}
\end{equation}
Indeed the determinant equals $\det D^2\psi$; the bracket is at most
$2-2p<0$.  On this cone,
$\|DV(b,r)\|\le C r^{p-2}$, with $C$ depending only on the fixed
parameters.  Consequently
\begin{equation}
 c\sup_{D_c}\|DV\|\le Cc d^{p-2}
   =C c^{1-\gamma(2-p)}\longrightarrow0,
 \qquad d^2=o(c^2).
 \label{c:sing-cutoff-estimates}
\end{equation}
Thus, for sufficiently small $c$, both $cDV$ and $-cDV$ have
eigenvalues $\lambda,-\lambda$ with $0<\lambda<1$ at every label.
Although no derivative bound is asserted at $r=0$, each fixed $D_c$
has a neighborhood on which the field is smooth.

Write $s=1-t$.  Consider the two sets of segments
\begin{align}
 E_c^-&=\{(x,y,\sqrt{2}\,t): |x|\le1,
            \ (y,t)=(b,r)+cxV(b,r),\ (b,r)\in D_c\},\notag\\
 E_c^+&=\{(x,y,\sqrt{2}\,t): |y|\le1,
            \ (x,s)=(b,r)-cyV(b,r),\ (b,r)\in D_c\}.
 \label{c:sing-families}
\end{align}
We apply \BScaling{} with $h=\sqrt{2}$ and $L=1$ to the first
family with field $cV$.
Interchanging $x,y$ and
reflecting and translating the physical height gives the second family
with field $-cV$.  Thus, for each fixed sufficiently small $c$ and any
$e>0$, each family has a finite cylinder cover with compact
parallelogram bases and total base area at most
\begin{equation}
 J(c)+e,
 \qquad
 J(c)=\sqrt{2}\int_{D_c}
       \frac{db\,dr}{\sqrt{1+c^2(P^2+2Q^2)}}.
 \label{c:sing-cost-integral}
\end{equation}
The factor $2$ multiplying $Q^2$ records the physical height scale.

These two applications use the companion's single square-tile
construction. For the fixed compact cutoff domain $D_c$, the fields
$\pm cV$ are $C^1$ on a neighborhood. Equations
\eqref{c:sing-hyperbolicity} and \eqref{c:sing-cutoff-estimates}
give $\operatorname{tr}D(\pm cV)=0$ and
$-1<\det D(\pm cV)<0$ for sufficiently small $c$, which is the
theorem's condition for $L=1$. The tile orientations may vary, but
their weights use the velocities in the original $(b,r)$ coordinates,
giving \eqref{c:sing-cost-integral} with the physical factor $\sqrt2$.
No condition on the area of $\partial D_c$ or smallness of the
operator norm is required by the theorem. The separate cylinders
introduced next cover the omitted tips.


We add an $x$-parallel cylinder with triangular base
$\{(0,y,\sqrt{2}\,t):0\le t\le2d,\ |y|\le t\}$,
and the analogous $y$-parallel cylinder at the upper end, obtained by
projecting $\{0\le s\le2d,\ |x|\le s\}$ to the plane $y=0$.
Each base has area
\begin{equation}
 \sqrt{2}\int_0^{2d}2r\,dr
       =\sqrt{2}(2d)^2.
 \label{c:sing-tip-cost}
\end{equation}

\subsection{Inversion, including the cone boundary}

All estimates below are uniform over the points under consideration,
with $p$ and $\gamma$ fixed.  Suppose
\[
 2d\le T\le1,\qquad |B|\le T,\qquad |k_0|\le c.
\]
We first solve
\begin{equation}
 (B,T)=(b,r)+k_0V(b,r)
 \label{c:sing-inverse}
\end{equation}
without imposing the last inequality of \eqref{c:sing-domain}.
In the closed Euclidean ball of radius $T/4$ about $(B,T)$, one has
$3T/4\le r\le5T/4$ and $|b|/r\le5/3$.  Formula
\eqref{c:sing-field} gives bounds
$|V|\le C T^{p-1}$ and $\|DV\|\le C T^{p-2}$ throughout the ball.
The map
\[
 (b,r)\longmapsto(B,T)-k_0V(b,r)
\]
maps this ball into itself and is a contraction for all sufficiently
small $c$, since $cT^{p-2}\le c(2d)^{p-2}=o(1)$.
Its fixed point satisfies $d\le r\le2$.

To prove that this fixed point lies in the cone, use the exact identities
\begin{equation}
 Q(b,r)\pm P(b,r)
   =\frac{r\pm b}{2}(r/a)^{p-2}
                      \left[(2-p)b/r\pm p\right].
 \label{c:sing-wall-factorization}
\end{equation}
Equation~\eqref{c:sing-inverse} consequently gives
\[
 T\pm B=(r\pm b)(1+\theta_\pm),\qquad
 |\theta_\pm|\le Cc r^{p-2}=o(1).
\]
The factors $1+\theta_\pm$ are positive.  Since $T\pm B\ge0$, we
obtain $r\pm b\ge0$.  This also handles equality on either wall:
no boundary point is lost in passing to inverse labels.

\subsection{The two families overlap sufficiently}

Take a point of $K_2$ outside the two tip regions, so that
$t,s\ge2d$.  Apply \eqref{c:sing-inverse} to $(B,T)=(y,t)$
with $k_0=cx$, and to $(B,T)=(x,s)$ with $k_0=-cy$.
Let the resulting labels be $(b_-,r_-)$ and $(b_+,r_+)$.
They already satisfy $d\le r_\pm\le2$ and $|b_\pm|\le r_\pm$.
It remains to show that one label satisfies its ceiling in
\eqref{c:sing-domain}.

Suppose both ceilings fail.  Since $\delta\ge0$, this implies
$r_-,r_+>a$.  The inverse equations give
\begin{align}
 b_-&=y-cxP(b_-,r_-),&
 b_+&=x+cyP(b_+,r_+),\notag\\
 r_-+r_+&=1+c\left[yQ(b_+,r_+)-xQ(b_-,r_-)\right].
 \label{c:sing-inverse-sum}
\end{align}
All labels now lie in a fixed compact region with $r\ge a$, on which
the field and its derivatives are bounded.  The last equation first
gives $r_-+r_+=1+O(c)$.  Because both summands exceed $a=1/2$,
each is $a+O(c)$.  The first two equations also give
$b_-=y+O(c)$ and $b_+=x+O(c)$.
Using $Q(b,a)=b$, the bracket in the last equation is therefore
$yx-xy+O(c)=O(c)$.  This improves the sum to $1+O(c^2)$ and
each height to $a+O(c^2)$.

Refining the first two equations and again using $Q(b,a)=b$, we obtain
\[
 b_-=y-cxP_a(y)+O(c^2),\qquad
 b_+=x+cyP_a(x)+O(c^2),\qquad
 Q(b_\pm,r_\pm)=b_\pm+O(c^2).
\]
Substitution into \eqref{c:sing-inverse-sum} yields
\begin{equation}
 r_-+r_+
   =1+c^2\left[x^2P_a(y)+y^2P_a(x)\right]+O(c^3).
 \label{c:sing-overlap-expansion}
\end{equation}
On the other hand, failure of both ceilings gives
\[
 r_-+r_+>
  1+c^2\left[\delta(b_-)+\delta(b_+)\right]+2c^{5/2}.
\]
Here $\delta(b_-)+\delta(b_+)=\delta(y)+\delta(x)+O(c)$, and
\begin{align*}
 \delta(y)+\delta(x)-x^2P_a(y)-y^2P_a(x)
  &=(x^2-y^2)\left[P_a(x)-P_a(y)\right]\\
  &=\frac{2-p}{2a}(x^2-y^2)^2\ge0.
\end{align*}
Comparison with \eqref{c:sing-overlap-expansion} would force
$2c^{5/2}\le Cc^3$, which is impossible for sufficiently small $c$.
Thus every middle-region point belongs to $E_c^-\cup E_c^+$.
Together with the two closed tip cylinders, the two families cover
all of $K_2$.

\subsection{The exact quadratic saving}

Let
\[
 H=\{(b,r):0<r\le a,\ |b|\le r\},\qquad |H|=a^2.
\]
The deletion of $r<d$ removes area $d^2=o(c^2)$.  The area added
above $r=a$ is
\[
 c^2\int_{-a}^{a}\delta(b)\,db+o(c^2).
\]
Indeed, for $|b|\le a$ its thickness is exactly
$c^2\delta(b)+c^{5/2}$.  Since $|b|\le r\le2$ and $\delta$ is
bounded there, the remaining portion has $a<|b|\le a+O(c^2)$
and thickness $O(c^2)$, hence area $O(c^4)$.
The field is uniformly bounded on $|b|\le r\le2$ because $p>1$.
Taylor expansion of the integrand in
\eqref{c:sing-cost-integral} therefore has uniform remainder
$O(c^4)$.  Also $|D_c\mathbin{\triangle}H|\to0$, so the bounded
quadratic coefficient may be integrated over $H$.  Including one
tip cylinder, we conclude that
\begin{equation}
 \frac{J(c)}{\sqrt{2}}+4d^2
  =a^2+c^2\kappa_p+o(c^2),\qquad
 \kappa_p=\int_{-a}^a\delta(b)\,db
                -\frac12\int_H(P^2+2Q^2)\,db\,dr.
 \label{c:sing-area-expansion}
\end{equation}

Both terms admit elementary evaluations.  The first is
\[
 \int_{-a}^a\delta(b)\,db
   =a^4\left(\frac p3+\frac{2-p}{5}\right)
   =\frac{2a^4(p+3)}{15}.
\]
For the second, put $\eta=b/r$.  Then
\[
 P=\frac a2(r/a)^{p-1}[p+(2-p)\eta^2],\qquad
 Q=a(r/a)^{p-1}\eta,\qquad db\,dr=r\,d\eta\,dr.
\]
It follows that
\begin{align*}
 \frac12\int_H(P^2+2Q^2)\,db\,dr
 &=\frac{a^4}{4p}\int_{-1}^{1}
          \left\{\frac14[p+(2-p)\eta^2]^2+2\eta^2\right\}\,d\eta\\
 &=\frac{a^4(2p^2+2p+13)}{30p}.
\end{align*}
Thus the exact coefficient is
\begin{equation}
 \boxed{\displaystyle
 \kappa_p=\frac{a^4(2p^2+10p-13)}{30p}<0.}
 \label{c:sing-coefficient}
\end{equation}
The range of $p$ in \eqref{c:sing-parameters} is precisely the
range above $1$ on which this expression is negative.  At $p=1$
its continuous extension is $-a^4/30$; the construction itself uses
$p>1$ to retain strict hyperbolicity.

\paragraph{Completion with finitely many cylinders.}
We can now finish the proof of Theorem~\ref{c:sing-theorem}.  First fix $p$ in the stated
range and then $\gamma$ in \eqref{c:sing-parameters}.
Choose a positive $c$ small enough for the derivative bounds, all
coverage statements, and
\[
 2J(c)+8\sqrt{2}\,d^2
  =\frac{\sqrt{2}}2+2\sqrt{2}\,\kappa_p c^2+o(c^2)
  <\frac{\sqrt{2}}2.
\]
Call the strict difference $\Delta>0$.
Only after fixing this $c$ do we apply
\BScaling{} to each family, choosing
its additive error less than $\Delta/4$.
Those two finite covers and the two triangular tip cylinders give
an actual finite cover of $K_2$ whose total base area is strictly less
than $A_{\min}(K_2)/2$.  The main bases are compact parallelograms.
No uniform choice of their meshes as $c\to0$ is needed.

\subsection{Edge-one normalization and alternative cutoffs}
\label{c:sing-edge-one}

The edge-one form allows a direct comparison with the rational field.
First we transport the whole construction by a similarity; then we explain
which cutoff and padding may be changed without changing its quadratic
saving.  The similarity halves physical lengths while leaving the height
label unchanged:
\[
 (x,y,\sqrt{2}\,t)\longmapsto(y/2,x/2,t/\sqrt{2}),
 \qquad (w,l)=(b/2,r),\qquad k=4c.
\]
Put $h=1/\sqrt{2}$ and define
\[
 \widehat V_1(w,l)=\tfrac14P(2w,l),\qquad
 \widehat V_2(w,l)=\tfrac12Q(2w,l),\qquad
 f(w)=w^2\widehat V_1(w,1/2).
\]
Equivalently, $\widehat V=(-\partial_lH_*,\partial_wH_*)$ for
$H_*(w,l)=2^{p-3}(w^2l^{p-2}-l^{p}/4)$.
The two segment families transform exactly, and
\[
 k^2(\widehat V_1^2+h^2\widehat V_2^2)
     =c^2(P^2+2Q^2),\qquad
 c^2\delta(2w)=k^2 f(w).
\]
With $S=\operatorname{diag}(1/2,1)$, the derivative identity is
\[
 \widehat V=\tfrac12 S V\circ S^{-1},\qquad
 \tfrac12 kD\widehat V=cS(DV)S^{-1}.
\]
Thus the scaled eigenvalues for axial half-length $L=1/2$ agree
with those of $cDV$.  Under the fixed physical rescaling, the derivative is likewise
similar to $\pm kD\widehat V$.
This is a spectral statement; the nonorthogonal matrix $S$ does not
preserve the ordinary operator norm. No operator-norm transfer is needed
by the square-tile theorem. The already constructed physical cylinder
cover also transports by the Euclidean similarity above.

Physical areas scale by $1/4$.  Therefore the per-family quadratic
coefficient after dividing physical area by $h$ is $\kappa_p/32$;
the positive saving in that normalization is
\[
 -\frac{\kappa_p}{32}
   =\frac{13-10p-2p^2}{15360p}.
\]
For this exact similarity the cutoff becomes $4^{-\gamma}k^\gamma$
and the ceiling padding becomes $k^{5/2}/32$.

The estimates also permit independent choices of these two auxiliaries.
Indeed, the preceding inverse, overlap, and area proofs hold with a
positive cutoff $\rho(c)$ and padding $q(c)>0$ whenever
\begin{equation}
 \rho(c)=o(c),\qquad c\rho(c)^{p-2}\longrightarrow0,
 \qquad q(c)=o(c^2),\qquad c^3=o(q(c)).
 \label{c:sing-cutoff-flexibility}
\end{equation}
The contraction and wall-sign estimates use only the second condition;
the deleted label area and two tip areas use the first.  The overlap
contradiction compares $2q(c)$ with $O(c^3)$, and the coefficient is
unchanged because $q(c)=o(c^2)$.  On the retained set $r\le2$, the
ceiling forces $r\le1/2+O(c^2)$.  Replacing its upper bound $2$ by
$1$ therefore leaves the label set unchanged for sufficiently small $c$.

We can now state the edge-one variant with its own cutoff.  Choose
$1<\beta<1/(2-p)$ and put $\sigma=k^\beta$.  The field and padding
polynomial are explicitly
\begin{align}
 \widehat V_1(w,l)&=\tfrac14(2l)^{p-1}
       [(2-p)(w/l)^2+p/4],\notag\\
 \widehat V_2(w,l)&=\tfrac12(2l)^{p-1}(w/l),\qquad
 f(w)=w^2[(2-p)w^2+p/16].
 \label{c:sing-edge-one-field}
\end{align}
The image tetrahedron is
\[
 \widetilde K_1=\{(x,y,hs):0\le s\le1,\ |x|\le s/2,
                         \ |y|\le(1-s)/2\}.
\]
It is the translate $K_1+(0,0,h/2)$ of the centered tetrahedron
in~\eqref{c:rat-tetrahedron}.  Its two segment families are
\[
 (x,s)=(w,l)+ky\widehat V(w,l),\quad |y|\le1/2,
 \qquad
 (y,1-s)=(w,l)-kx\widehat V(w,l),\quad |x|\le1/2,
\]
with labels
\[
 D_k=\{\sigma\le l\le1,\ |w|\le l/2,
             \ l\le1/2+k^2f(w)+k^{5/2}\}.
\]
Under $c=k/4$, these new choices pull back to
$\rho(c)=(4c)^\beta$ and $q(c)=(4c)^{5/2}$, which satisfy
\eqref{c:sing-cutoff-flexibility}.  Thus coverage follows from the
full preceding proof, including its closed inversions and two tip
cylinders.  The scaled eigenvalue bound is now
$O(k\sigma^{p-2})=o(1)$, so \BScaling{} applies with $L=1/2$.
For every positive error $e$, each finite family costs at most $I(k)+e$,
where
\begin{equation}
 I(k)=h\int_{D_k}
       (1+k^2(\widehat V_1^2+h^2\widehat V_2^2))^{-1/2}\,dw\,dl.
 \label{c:sing-edge-one-cost}
\end{equation}
There are exactly two additional tip cylinders, each of area
$2h\sigma^2$.

The coefficient has the following direct edge-one evaluation.  On
$D_0=\{0\le l\le1/2,\ |w|\le l/2\}$, put
\[
 A(p)=\int_{-1/4}^{1/4}f(w)\,dw
       =\frac{2-p}{2560}+\frac{p}{1536},
 \qquad
 G(p)=\frac12\int_{D_0}(\widehat V_1^2+h^2\widehat V_2^2)\,dw\,dl.
\]
Substituting $R=w/l$ gives
\begin{align}
 G(p)&=\frac1{16p}
 \left\{\frac{(2-p)^2/80+p(2-p)/24+p^2/16}{16}
                    +\frac1{96}\right\},\notag\\
 A(p)-G(p)&=\frac{2p^2+10p-13}{15360p}.
 \label{c:sing-edge-one-coefficient}
\end{align}
Here $\int_0^{1/2}l(2l)^{2p-2}\,dl=1/(8p)$ and the second
and fourth moments on $[-1/2,1/2]$ are $1/12$ and $1/80$.
In particular $A(1)=16/15360$ and $G(1)=17/15360$ are continuous
coefficient limits. They coincide with the rational limits of $J_d$
and $I_d$ in~\eqref{c:rat-coefficient-limits}.
We still fix $p>1$, then $\beta$, then a positive $k$, and finally
the two finite errors.  This gives $2I(k)+4h\sigma^2<h/4$ before
those errors.  The minimum projection area of $\widetilde K_1$ is
$h/2$ by the similarity above.

At the endpoint $p=1$, writing $\theta=b/r$, the original edge-two
field in \eqref{c:sing-field} becomes
\[
 P(b,r)=\frac{1+\theta^2}{4},\qquad Q(b,r)=\frac{\theta}{2}.
\]
These formulas depend only on the angle of the label.  The cutoff
requirements \eqref{c:sing-cutoff-flexibility} would now demand both
$\rho(c)=o(c)$ and $c/\rho(c)\to0$, which are incompatible.
The next section fixes a truncated label set, applies the square-zero
case of the same finite-approximation theorem, and covers the apex
segments by a finite velocity net. The field agrees with the endpoint
above; the padding and the treatment of the apex are chosen anew.
